\documentclass[11pt, letterpaper]{article}
\input{../../homework.tex}
\usepackage{titlepageBU}
\title{Problem Set 5}
\courseID{PY 451}
\courseSection{A1}
\begin{document}
\makeproblem
\section*{Problem 5}
\begin{enumerate}
	\item 
\[
	\left( S_x^{(1)} \otimes S_y^{(2)}  \right) \left( \ket{+}_1\otimes \ket{+}_2 \right) =\left( S_x^{(1)} \ket{+}_1 \right) \otimes \left(S_y^{(2)} \ket{+}_2\right)=\frac{\hbar}{2} \ket{-}_1 \otimes \frac{i\hbar}{2} \ket{-}_2
.\]
\item The norm squared is
	\[
		\left( \frac{\hbar }{2} \bra{-}_1  \otimes \frac{-i\hbar }{2} \bra{-}_2   \right) \left( \frac{\hbar }{2} \ket{-}_1 \otimes \frac{i\hbar }{2} \ket{-}_2  \right) =\frac{\hbar ^2}{4} \braket{-}{-}_1 \cdot \frac{\hbar ^2}{4} \braket{-}{-}_2 =\frac{\hbar ^4}{16} 
	.\]
	Therefore, the norm is $\frac{\hbar ^2}{4} $.
\end{enumerate}
\section*{Problem 2}
\begin{enumerate}
	\item Note that $P_+^{(1)} =\ket{+}_1\,_1 \bra{1}$. Moreover, factoring gives $\ket{\psi } =\frac{1}{\sqrt{2} } \left( \ket{+}_1+\ket{-}_1 \right) \otimes \ket{+}_2$.
		\begin{align*}
			\mathcal{P}_{+} &=\bra{\psi } P_+^{(1)} \otimes \mathbbm{1}^{(2)} \ket{\psi } \\
					&=\bra{\psi } \left( P_+^{(1)}\ket{\psi }_1 \otimes \mathbbm{1}^{(2)} \ket{\psi }_2 \right) \\
					&=\left(_1 \bra{\psi }P_+^{(1)} \ket{\psi }_1\right)\left(_2\braket{\psi }{\psi }_2 \right) \\
					&=\,_1\braket{\psi }{+} _1\,_1\braket{+}{\psi }_1\\
					&=\left| _1 \braket{\psi }{+}_1 \right| ^2\\
					&=\left| \left(\frac{1}{\sqrt{2} }  \,_1\bra{+} +\,_1\bra{-}  \right) \ket{+}_1 \right| ^2\\
					&=\left| \frac{1}{\sqrt{2} }\,_1\braket{+}{+}_1  \right|^2\\
					&=\frac{1}{2} 
		.\end{align*}
		The normalized state is then $\ket{+}_1\otimes \ket{+}_2$, since we have collapsed the first election into $\ket{+} $ by measuring it, and have not measured the second electron at all, thus not changing it's state. Since the only other possibility is measuring $-\hbar /2$, and probabilities sum to 1, we have that $\mathcal{P}_{-} =1/2$ as well, and by the same logic that normalized state post-measurement would be $\ket{-}_1\otimes \ket{+} _2$.

	\item Note that
		\[
			\frac{1}{2} \left( \ket{+}_1+\ket{-}_1 \right) \otimes \left( \ket{+}_2 +\ket{-}_2 \right) =\frac{1}{\sqrt{2} } \left( \ket{+}_1+\ket{-}_1 \right) \otimes \frac{1}{\sqrt{2} } \left( \ket{+}_2+\ket{-}_2 \right) 
		.\]
		This follows from bilinearity of the tensor product and I will use it to justify the fact that $\,_2\braket{\psi }{\psi }_2 =1$, since this makes both individual electron states clearly normalized.
		\begin{align*}
			\mathcal{P}_+&=\bra{\psi } \left( P_+^{(1)} \otimes \mathbbm{1}^{(2)}  \right) \ket{\psi } \\
				     &=\bra{\psi } \left( P_+^{(1)} \ket{\psi }_1+\mathbbm{1}^{(2)} \ket{\psi }_2 \right) \\
				     &=\,_1\bra{\psi }P_+^{(1)}\ket{\psi }_1\,_2\bra{\psi }\mathbbm{1}^{(2)} \ket{\psi }_2\\
				     &=\,_1\braket{\psi }{+}_1\,_1\braket{+}{\psi }_1\,_2\braket{\psi }{\psi }_2\\
				     &=\left|_1 \braket{+}{\psi }_1 \right|^2\\
				     &=\left| \frac{1}{\sqrt{2} }\,_1 \braket{+}{+}_1  \right|^2\\
				     &=\frac{1}{2} 
		.\end{align*}
		Once again, the normalized state post-measurement is simply $\ket{+}_1 \otimes \frac{1}{\sqrt{2} } \left( \ket{+}_2+\ket{-}_2 \right) $, as we have collapsed the first elecron to $\ket{+} $, and by my earlier restatement of $\ket{\psi } $, the resultant state from doing this is still normalized. Since probabilities sum to 1, $\mathcal{P}_-=\frac{1}{2} $ as well, and the collapsed state after that measurement would be $\ket{-}_1 \otimes \frac{1}{\sqrt{2} } \left( \ket{+}_2+\ket{-}_2 \right) $ by the same logic.

	\item This one is just going to be messy. I will omit the subscripts for now because I just don't want to type them.
		\begin{align*}
			\bra{\psi } P_+\otimes \mathbbm{1}\ket{\psi } &=\bra{\psi } \left( P_+ \otimes \mathbbm{1} \right) \left( \frac{1}{\sqrt{2} } \left( \ket{+} \otimes \ket{+} \right) +\left( \ket{-} \otimes \ket{-}  \right)  \right) \\
								      &=\frac{1}{\sqrt{2} } \bra{\psi } \left( \left( P_+\otimes \mathbbm{1} \right) \left( \ket{+} \otimes \ket{+}  \right) + \left( P_+\otimes \mathbbm{1} \right) \left( \ket{-} \otimes \ket{-}  \right)  \right) \\
								      &=\frac{1}{\sqrt{2} } \bra{\psi } \left( \left(P_+\ket{+} \otimes \mathbbm{1}\ket{+}  \right)+\left( P_+\ket{-} \otimes \mathbbm{1}\ket{-}  \right)  \right)\\
								      &=\frac{1}{\sqrt{2} } \bra{\psi } \left(\left( \ket{+} \otimes \ket{+}\right) +\left( 0\otimes \ket{-}  \right)  \right) \\
								      &=\frac{1}{2} \left( \left(\bra{+} \otimes \bra{+} \right)+\left( \bra{-} \otimes \bra{-}  \right)  \right) \left( \left( \ket{+} \otimes \ket{+}  \right) +\left( 0\otimes \ket{-}  \right)  \right) \\
								      &=\frac{1}{2} ( \left( \bra{+} \otimes \bra{+}  \right) \left( \ket{+} \otimes \ket{+}  \right) +\left( \bra{+} \otimes \bra{+}  \right) \left( 0\otimes \ket{-}  \right) \\
								      & +\left( \bra{-} \otimes \bra{-}  \right) \left( \ket{+} \otimes \ket{+}  \right) +\left( \bra{-} \otimes \bra{-}  \right) \left( 0\otimes \ket{-}  \right)  )\\
								      &=\frac{1}{2} \left( \left| \braket{+}{+}  \right|^2+\braket{+}{0} \braket{+}{-} +\left| \braket{-}{+}  \right| ^2+\braket{-}{0} \braket{-}{+}  \right) \\
								      &=\frac{1}{2} 
		.\end{align*}
		Once again, since probabilities sum to $1/2$, the probability that we measure $-\hbar /2$ in this specific state must also be $\mathcal{P}_{-} =1/2$. If we measure $\ket{+}_1$, then the only state in the superposition that we could measure is $\ket{+}_1 \otimes \ket{+}_2$. This can also be seen by noting that in my calculation of $\mathcal{P} _+$, I found that $\left( P_+^{(1)} \otimes \mathbbm{1}^{(2)}  \right) \ket{\psi } =\left( \ket{+} \otimes \ket{+}  \right) +\left( 0\otimes \ket{-}  \right) $. Since the tensor product is bilinear, the product $-\otimes \ket{-} $ must be linear, and linear maps send $0\mapsto 0$. Therefore, plugging in the zero vector gives $0\otimes \ket{-} =0$. We have thus found that measuring $\hbar /2$ in the first electron, given by the projection operator $P_+^{(1)} $, maps
		\[\begin{tikzcd}[row sep = 2.5em, column sep = 2.5em]
			\ket{\psi } \arrow[r, "P_+"]&\ket{+}_1 \otimes \ket{+}_2 
		\end{tikzcd}\]

		Similar logic and calculations show that measuring $-\hbar /2$ yields a collapsed normalized state of $\ket{-}_1 \otimes \ket{-}_2$. I will omit the calculation since it is basically the same and needlessly messy.



	\item The calculations here are all basically the same as before, so I will start with the table, and will do the calculations afterwards. I will show less work because once again, they are basically the same as before.

		\hspace{-5em}\begin{tabular}{c|c|c|c|c}
			\toprule
			&$\mathcal{P} \left( S_z^{(1)} =\frac{\hbar }{2} ,S_z^{(2)} =\frac{\hbar }{2}  \right) $&$\mathcal{P} \left( S_z^{(1)} =-\frac{\hbar }{2} ,S_z^{(2)} =\frac{\hbar }{2}  \right) $&$\mathcal{P} \left( S_z^{(1)} =\frac{\hbar }{2} ,S_z^{(2)} =-\frac{\hbar }{2}  \right) $&$\mathcal{P} \left( S_z^{(1)} =-\frac{\hbar }{2} ,S_z^{(2)} =-\frac{\hbar }{2}  \right) $\\
			\hline
			i&1/2&1/2&0&0\\
			ii&1/4&1/4&1/4&1/4\\
			ii&1/2&0&0&1/2\\
			\bottomrule
		\end{tabular}
		\begin{enumerate}
			\item Recall the first normalized state after measurement was $\ket{\psi } =\ket{\pm}_1 \otimes \ket{+}_2$. Measuring again in $S_z^{(2)} $ gives a probability of seeing $\hbar /2$ as
				\begin{align*}
					\mathcal{P}_{\pm,+} &=\bra{\psi } \left( \mathbbm{1}^{(1)} \otimes \ket{+} \bra{+}^{(2)}  \right) \ket{\psi } \\
							  &=\left| _2\braket{\psi }{+}_2 \right| ^2\\
							  &=\left|_2 \braket{+}{+}_2 \right| ^2\\
							  &=1
				.\end{align*}
				Multiplying this by the probability of already having measured a state $\ket{+}_1 $ or $\ket{-}_1$ gives probabilities $\mathcal{P}_{+,+}=\mathcal{P}_{-,+} =1/2$. We also find that the probability of measuring $\ket{-} $ is $0$, since probabilities sum to 1.
		\item The normalized post-measurement states were $\ket{\pm}_1 \otimes \frac{1}{\sqrt{2} } (\ket{+}_2 +\ket{-}_2) =\ket{\pm} \otimes \ket{+_x}_2$. Measuring again in $S_z^{(2)} $ gives
				\begin{align*}
					\mathcal{P}_{\pm,+} &=\bra{\psi } \left( \mathbbm{1}^{(1)} \otimes \ket{+} \bra{+}^{(2)}  \right) \ket{\psi } \\
							    &=\left| _2\braket{+_x}{+}_2 \right| ^2\\
							    &=\frac{1}{2} 
				.\end{align*}
				Multiplying through by the probability of measuring either $\ket{+}_1$ or $\ket{-}_1$ gives probabilities of $\mathcal{P} _{+,+} =\mathcal{P}_{-,+} =1/4$. Since probabilities sum to 1, we have $\mathcal{P}_{\pm,-}=1/2$, and by the same logic we have all other probabilities are $1/4$ as well.
			\item The state post measurement is either $\ket{+} \otimes \ket{+} $ or $\ket{-} \otimes \ket{-}$. Note that
				\begin{align*}
					\mathcal{P}_{++} &=\bra{\psi } \left( \mathbbm{1}^{(1)} \otimes \ket{+} \bra{+}^{(2)}  \right) \ket{\psi } \\
					&=\left|_2 \braket{+}{+} _2 \right| ^2\\
					&=1
				.\end{align*}
				
				Multiplying through by the probability of finding $\ket{+}_1 \otimes \ket{+}_2$ in the first place gives a probability of $1/2$. Similarly,
				\begin{align*}
					\mathcal{P}_{- -} &=\bra{\psi } \left( \mathbbm{1}^{(1)} \otimes \ket{-} \bra{-} ^{(2)}  \right) \ket{\psi } \\
							  &=\left| _2 \braket{-}{-} _2 \right| ^2\\
							  &=1
				.\end{align*}
				By the same logic, the true probability here is $1/2$. Since probabilities sum to $0$, we have $\mathcal{P}_{+-} =\mathcal{P}_{-+} =0$.
		\end{enumerate}
\end{enumerate}
\end{document}
